%=====================================================================
\chapter{Case Study Research}\label{ch:casestudy}
%=====================================================================
\locator{3}{Part~3 --- Research Designs for Computing}

\begin{outcomes}
By the end of this chapter you will be able to:
\begin{tight}
  \item \textbf{Distinguish} a case study from an example, a pilot and a small
    survey.
  \item \textbf{Define} the unit of analysis and choose between holistic and
    embedded designs.
  \item \textbf{Select} cases on theoretical grounds and defend the selection.
  \item \textbf{Maintain} a chain of evidence from research question to
    conclusion.
  \item \textbf{Argue} analytic generalisation without claiming statistical
    generalisation.
\end{tight}
\end{outcomes}

\begin{prereq}
\chref{philosophy} --- case study research sits most comfortably in critical
realism or interpretivism, and the difference changes what you claim.
\chref{questions} for how-and-why questions.
\end{prereq}

\begin{keyterms}
case study $\bullet$ unit of analysis $\bullet$ holistic and embedded design
$\bullet$ single and multiple case $\bullet$ proposition $\bullet$ theoretical
replication $\bullet$ literal replication $\bullet$ chain of evidence $\bullet$
triangulation $\bullet$ pattern matching $\bullet$ analytic generalisation
$\bullet$ case study database
\end{keyterms}

\section{The Most Abused Term in Computing Research}

``We evaluated our approach with a case study'' usually means ``we tried it
once''. That is a demonstration (\chref{designscience}), not a case study, and
the confusion costs authors credibility with anyone who knows the difference.

\begin{definitionbox}[Case study]
An empirical investigation of a contemporary phenomenon \emph{in its real-life
context}, especially where the boundary between phenomenon and context is not
clear, using multiple sources of evidence.

\medskip
The italicised clause is the point. If you can separate the phenomenon from its
context, you should --- that is an experiment, and it will give you a cleaner
answer. The case study exists for situations where the context is part of what
you are studying, so removing it would destroy the phenomenon.
\end{definitionbox}

\begin{center}
\small
\begin{tabular}{@{}L{3.2cm}L{5.4cm}L{5.5cm}@{}}
\toprule
\textbf{It is not} & \textbf{Because} & \textbf{Which would be} \\
\midrule
An example
  & One illustration of your method working
  & A demonstration in a DSR evaluation \\
A pilot
  & Trying something once to see what happens
  & A feasibility study, honestly labelled \\
A small survey
  & Asking a few people questions
  & An underpowered survey (\chref{survey}) \\
An experiment with N=1
  & No manipulation, no control, no randomisation
  & Nothing --- N=1 experiments require special designs \\
\bottomrule
\end{tabular}
\end{center}

\begin{paperbox}[The reference for our field]
Per Runeson and Martin H{\"o}st, \emph{Guidelines for conducting and reporting
case study research in software engineering}, Empirical Software Engineering
14(2), 2009~\cite{runeson2009guidelines}.\oa{}

\medskip
\textbf{Why it matters.} It translates Yin's general methodology into software
engineering, with a reporting structure and a checklist. It is open access and
it is what reviewers in our field expect you to have followed.

\medskip
\textbf{What to extract.} The five process steps, the distinction between the
four purposes (exploratory, descriptive, explanatory, improving), and the
reporting checklist --- which you should use as a table of contents for the
methodology chapter.
\end{paperbox}

\section{Design: Two Choices That Determine Everything}

\subsection{The unit of analysis}

The \term{unit of analysis} is what the case \emph{is}. It is astonishing how
often this is left implicit, and every downstream decision depends on it.

\begin{center}
\small
\begin{tabular}{@{}L{5.2cm}L{9.1cm}@{}}
\toprule
\textbf{If the question is} & \textbf{The case is probably} \\
\midrule
How does this team adopt the practice?         & The team \\
How does this organisation manage the migration? & The organisation \\
How does this project's architecture degrade?  & The project \\
How do individual reviewers decide?            & The reviewer, embedded in a
                                                 team-level case \\
\bottomrule
\end{tabular}
\end{center}

\subsection{Four designs}

\begin{center}
\small
\begin{tabular}{@{}L{2.6cm}L{5.6cm}L{5.6cm}@{}}
\toprule
 & \textbf{Holistic} (one unit) & \textbf{Embedded} (sub-units within) \\
\midrule
\textbf{Single case}
  & One organisation studied as a whole. Justified when the case is critical,
    unique, revelatory, or longitudinal
  & One organisation, with several teams analysed within it. Guards against
    the whole-case account floating free of evidence \\
\textbf{Multiple case}
  & Several organisations compared as wholes
  & Several organisations, each with sub-units. Richest, most expensive, and
    the strongest basis for cross-case conclusions \\
\bottomrule
\end{tabular}
\end{center}

\begin{alertbox}[A single case needs an explicit justification]
Single-case designs are legitimate, and Yin gives five grounds: the case is
\emph{critical} (it tests a well-formed theory), \emph{extreme or unusual},
\emph{common} (and therefore informative about ordinary situations),
\emph{revelatory} (previously inaccessible), or \emph{longitudinal}.

\medskip
``It was the only company that would let me in'' is not on that list. If access
drove the choice --- as it often, honestly, does --- then find which of the five
grounds \emph{also} applies and argue that, or reframe the study as exploratory.
Do not leave the reader to infer that the case was chosen by convenience, since
they will.
\end{alertbox}

\section{Propositions and the Chain of Evidence}

\term{Propositions} are what you expect to find and why. They are not
hypotheses to be tested statistically; they direct attention, and they make the
study falsifiable in the sense that matters here --- you can be surprised.

\begin{conceptbox}[Exploratory case studies have questions, not propositions]
If the study is genuinely exploratory, say so and state what you will look at
instead of what you expect. What you must not do is write propositions after
data collection and present them as though they guided it. That is HARKing
(\chref{prereg}) in qualitative clothing.
\end{conceptbox}

The \term{chain of evidence} is the property that a reader can follow a
conclusion back through your report to the evidence, and forward from a
research question to the data that addressed it, in both directions without a
gap.

\begin{center}
\small
\begin{tabular}{@{}L{0.6cm}L{13.8cm}@{}}
\toprule
1 & The report cites specific evidence --- ``Interview 7, T3'' --- not
    ``participants said''. \\
2 & That evidence exists in a \term{case study database}: transcripts, notes,
    documents, logs, kept separately from your narrative. \\
3 & The database records the circumstances of collection: when, where, who,
    under what protocol. \\
4 & The protocol links the collection back to the research questions. \\
\bottomrule
\end{tabular}
\end{center}

\begin{alertbox}[The case study database is what makes this auditable]
Keeping the raw evidence separate from your interpretation of it is the
qualitative equivalent of publishing your data (\chref{reproducibility}). It is
also what \reg{PhD Regs 2024}{\S14(i)} requires when it states that ``the raw
data will have to be submitted to the supervisor along with a copy of the
thesis'', and it is what lets a supervisor or examiner check that a conclusion
is actually supported.
\end{alertbox}

\section{Selecting Cases, and Multiple-Case Logic}

Cases are chosen \emph{theoretically}, never randomly. A random sample of four
organisations is neither representative --- four is far too few --- nor
informative, because nothing connects them to your question.

\begin{center}
\small
\begin{tabular}{@{}L{3.6cm}L{10.7cm}@{}}
\toprule
\textbf{Logic} & \textbf{Choose cases so that} \\
\midrule
Literal replication
  & They are similar in the respects your theory says matter, and you predict
    similar results. Confirmation across them strengthens the theory \\
Theoretical replication
  & They differ in a respect your theory says matters, and you predict
    \emph{different} results for a stated reason. This is far stronger: it is
    a genuine test \\
\bottomrule
\end{tabular}
\end{center}

\begin{conceptbox}[Theoretical replication is the case study's version of an experiment]
If your theory says the practice will help co-located teams and hinder
distributed ones, then studying one of each and finding exactly that is
powerful evidence --- and finding the opposite refutes the theory. Predicting
the difference \emph{in advance} is what converts a descriptive study into a
test, and it costs nothing but the discipline of writing the prediction down
before collecting.
\end{conceptbox}

\section{Analysis}

\begin{center}
\small
\begin{tabular}{@{}L{3.4cm}L{10.9cm}@{}}
\toprule
Pattern matching
  & Compare an empirically observed pattern with the one your propositions
    predicted. Correspondence supports the proposition; divergence is a
    finding \\
Explanation building
  & Iteratively build an explanation of the case, revising as evidence
    accumulates. Powerful and dangerous --- record the revisions, or it becomes
    unfalsifiable \\
Time-series / process
  & Trace the sequence of events, testing a predicted process against the
    observed one \\
Cross-case synthesis
  & Treat each case as a whole and compare. Word tables per case, then compare
    across \\
Rival explanations
  & For each conclusion, state the best competing explanation and the evidence
    that discriminates. This is the single most persuasive thing in a case
    study report \\
\bottomrule
\end{tabular}
\end{center}

\subsection{Triangulation}

\begin{center}
\small
\begin{tabular}{@{}L{3.0cm}L{11.3cm}@{}}
\toprule
Data triangulation        & Multiple sources: interviews, repository history,
                            documents, observation \\
Observer triangulation    & More than one researcher collecting or coding \\
Methodological triangulation & Qualitative and quantitative strands
                            (\chref{mixed}) \\
Theory triangulation      & Interpreting through more than one theoretical lens \\
\bottomrule
\end{tabular}
\end{center}

\begin{pitfall}[Triangulation is not corroboration-hunting]
The purpose is not to accumulate sources that agree. It is to expose
disagreement, because disagreement is informative: when developers say review
is thorough and the repository shows a median review time of ninety seconds,
you have found something worth the whole study.

\medskip
A case study reporting that all sources agreed on everything has usually not
looked hard enough, or has looked only where agreement was likely.
\end{pitfall}

\section{What You May Claim}

\begin{center}
\small
\begin{tabular}{@{}L{4.2cm}L{10.1cm}@{}}
\toprule
\textbf{Statistical generalisation} & From a sample to a population, by
                                      probability. \textbf{Case studies cannot
                                      do this}, and should not pretend to \\
\textbf{Analytic generalisation}    & From findings to a theory, which is then
                                      applicable to other cases sharing the
                                      relevant conditions. This is what a case
                                      study offers \\
\bottomrule
\end{tabular}
\end{center}

\begin{conceptbox}[How to write the generalisation paragraph]
Not: ``these findings apply to software companies in Pakistan''. That is a
statistical claim four cases cannot support.

\medskip
Instead: ``the mechanism identified here --- that review checklists are
abandoned when review is time-boxed below the checklist's completion time ---
should be expected wherever those two conditions co-occur. We observed it in
four organisations differing in size and domain, which increases confidence
that it is not an artefact of any one setting. It would not be expected where
review is untimed, and testing that is the natural next study.''

\medskip
That paragraph states a mechanism, its scope conditions, the evidence for
transferability, and a falsifiable prediction. It is stronger than the
statistical claim would have been, and it is honest.
\end{conceptbox}

\begin{traditions}{the case study}
\tradEMP{Uneasy: no randomisation, no statistical inference. Its contribution
is to generate the theory an experiment later tests, and to establish that a
phenomenon occurs in the field at all.}
\tradDSR{Provides the naturalistic summative evaluation of the FEDS grid
(\chref{designscience}) --- what happened when the artefact met real use.}
\tradQUAL{Its home. Thick description, multiple sources, context as
constitutive, and the researcher acknowledged as part of the instrument.}
\tradFORM{Contributes little directly, but case studies are where the
assumptions of formal models are checked against practice --- often
unfavourably.}
\end{traditions}

\begin{lab}[Design a two-case study with theoretical replication]
\begin{tightnum}
  \item State the research question as a how or why question.
  \item Name the unit of analysis in one noun phrase.
  \item Write two propositions.
  \item Choose two cases for \emph{theoretical} replication: name the dimension
    on which they differ and predict, in writing, how the results will differ.
  \item List at least three data sources per case and say which proposition
    each addresses.
  \item For your main expected finding, write the strongest rival explanation
    and the evidence that would discriminate between it and yours.
\end{tightnum}

\medskip
Step 6 is the one examiners notice. Feed this into portfolio piece P3.
\end{lab}

\begin{checkpoint}
\begin{tightnum}
  \item Why is a case study preferable to an experiment when phenomenon and
    context cannot be separated? What is given up?
  \item A study of four companies concludes that ``agile adoption in Pakistan
    faces cultural barriers''. What is wrong with the claim, and how would you
    rewrite it?
  \item Explain why theoretical replication provides stronger evidence than
    literal replication.
  \item Your interview data and repository data disagree. What are the three
    things this could mean, and how would you tell them apart?
\end{tightnum}
\end{checkpoint}

\begin{chaptersummary}
\begin{tight}
  \item A case study investigates a phenomenon in a context that cannot be
    separated from it. Trying something once is a demonstration, not a case
    study.
  \item Name the unit of analysis explicitly; choose holistic or embedded,
    single or multiple.
  \item A single case needs one of the five justifications. Access alone is not
    one.
  \item Propositions direct attention and let you be surprised. Exploratory
    studies state questions instead --- and never invent propositions
    afterwards.
  \item The chain of evidence and the case study database make the work
    auditable, and \S14(i) requires the raw data anyway.
  \item Choose cases theoretically. Predicting a difference in advance ---
    theoretical replication --- turns description into a test.
  \item Triangulate to find disagreement, not to accumulate agreement.
  \item Claim analytic generalisation: a mechanism with scope conditions and a
    falsifiable prediction. Never statistical generalisation.
\end{tight}
\end{chaptersummary}

\begin{reviewq}
  \item Case study research is criticised for lacking rigour and permitting
    biased interpretation. Identify the specific methodological devices in this
    chapter that answer that criticism, and say which of them is hardest to
    verify from a published report.
  \item Explanation building revises the explanation as evidence accumulates,
    which resembles the practice \chref{prereg} condemns as HARKing in
    quantitative work. Are they the same thing? Distinguish carefully.
  \item You have access to one organisation and a strong theory. Design a
    single-case study that constitutes a genuine test of that theory, and state
    what result would refute it.
  \item \reg{PhD Regs 2024}{\S14.3(A)} asks whether ``the sample produce[s] the
    type of knowledge necessary to understand the structures and processes
    within which the individuals or situations are located''. Contrast this
    criterion with statistical representativeness, and say why applying the
    wrong one to a case study is a category error.
\end{reviewq}
